% data_availability.tex --- Section 6
\section{Data Availability, Funding, and Acknowledgments}
\label{sec:data}

\subsection{Open-source artifacts}

All artifacts are publicly available to support replication:

\begin{itemize}
  \item \textbf{Production model:} \texttt{PeetPedro/kompress-v8} on
        HuggingFace---149M-parameter dual-head ModernBERT, the deployed
        production checkpoint (heretic $0.955$, $\mathit{override\_delta}{=}0$).
  \item \textbf{Training \& orchestration engine:}
        \texttt{peterlodri-sec/loopkit} on GitHub---four-phase state machine,
        SQLite kernel, LLM Council operator, Telegram bot.
  \item \textbf{Dataset:} \texttt{PeetPedro/ultrawhale-dogfood} on
        HuggingFace---structural silver-label distributions with regex-driven
        must-keep thresholds, generated by the ultrawhale dogfeed loop.
  \item \textbf{Outer-loop system:} \texttt{peterlodri-sec/ultrawhale} on
        GitHub---the self-building coding agent whose dogfeed loop produced
        the dataset and the research program that led to this paper.
  \item \textbf{Log registry \& replication vault:} chronological experiment
        logs, state specifications, and telemetry at
        \url{https://pocoo.vaked.dev}.
\end{itemize}

\subsection{Compute costs}

The total cost of producing this research is bifurcated across two layers:

\begin{itemize}
  \item \textbf{Outer loop (ultrawhale v1.0.0$\to$v100.1.0):} \$37.19 in
        DeepSeek API fees---the coding agent that built itself, generated the
        dataset, and produced the research program. Approximately \$0.24 per
        release across 150 blocks.
  \item \textbf{Inner loop (kompress fine-tuning):} \$0.55 in vast.ai GPU
        compute (RTX~4090 at \$0.38/hr) for the v3$\to$v4$\to$v5
        self-labeling chain that constitutes the paradox evidence.
\end{itemize}

The combined cost of under \$38 demonstrates that the Loop Engineering
paradigm can produce publication-grade ML research at a total token and
compute spend below that of a single industry conference registration.

\subsection{Decentralized open-source research funding}

This work was not supported by traditional institutional grants. We adopt a
decentralized open-source funding model in which a community council guides
donations to sustain infrastructure, compute, and maintainer time. This model
aligns with the open-source ethos of the artifact ecosystem above and avoids
the incentive distortions of grant-driven milestone scheduling. We thank the
community council and the donors who made the compute runs possible.

\subsection{Acknowledgments}
% TODO: named contributors, the community council, compute donors,
% Qwen2.5 / ModernBERT / LangChain / Osmani / Greyling ecosystem acknowledgments.
